\documentclass[12pt]{article}

% Layout.
\usepackage[top=1in, bottom=0.75in, left=0.75in, right=0.75in, headheight=1in, headsep=6pt]{geometry}

% Fonts.
\usepackage{mathptmx}

\usepackage[scaled=0.86]{helvet}
\renewcommand{\emph}[1]{\textsf{\textbf{#1}}}

% Misc packages.
\usepackage{amsmath,amssymb,latexsym,amsthm}
\usepackage{graphicx,hyperref}
\usepackage{array}
\usepackage{xcolor}
\usepackage{multicol}
\usepackage{tabularx,colortbl}
\usepackage{enumitem}
\usepackage{soul,multicol}
\usepackage{setspace}
\doublespacing

\hypersetup{
    colorlinks=true,
    linkcolor=blue,
    filecolor=magenta,      
    urlcolor=blue,
    pdftitle={Proofs Worksheet},
    pdfpagemode=FullScreen,
    }

\def\mailto#1{\href{mailto:#1}{#1}}

%% special math symbols
\newcommand{\N}{\mathbb{N}}
\newcommand{\Z}{\mathbb{Z}}
\newcommand{\R}{\mathbb{R}}
\newcommand{\Q}{\mathbb{Q}}
\newcommand{\sse}{\subseteq}
\newcommand{\mcp}{\mathcal{P}}
\newcommand{\ol}[1]{\overline{#1}}

%other specials
\newcommand{\be}{\begin{enumerate}}
\newcommand{\ee}{\end{enumerate}}
\newcommand{\se}{\subseteq}
\newcommand{\spe}{\supseteq}
\newcommand{\ds}{\displaystyle}
\newcommand{\lcm}{\text{lcm}}
%\newcommand{\gcd}{\text{gcd}}

% Paragraph spacing
\parindent 0pt
\parskip 6pt plus 1pt
\def\tableindent{\hskip 0.5 in}
\def\ts{\hskip 1.5 em}

%header
\usepackage{fancyhdr}
\pagestyle{fancy} 
\lhead{\large\sf\textbf{MATH 265: Introduction to Mathematical Proofs}}
\rhead{\large\sf\textbf{Worksheet 18: Ch 10}}

\newcommand{\localhead}[1]{\par\smallskip\textbf{#1}\nobreak\\}%
\def\heading#1{\localhead{\large\emph{#1}}}
\def\subheading#1{\localhead{\emph{#1}}}

\newenvironment{clist}%
{\bgroup\parskip 0pt\begin{list}{$\bullet$}{\partopsep 4pt\topsep 0pt\itemsep -2pt}}%
{\end{list}\egroup}%

\begin{document}
\begin{enumerate}
\item Prove the statement below using induction. 
\begin{quote} The expression $n^3+5n+6$ is divisible by 3 for every $n \in \N.$\end{quote}



\vfill
\vfill
\item Can you think of any \emph{other} ways one might prove the statement above?\\

\vspace{.5in}
\item Proposition: For every $n \in \N,$ every set of $n$ horses has the same color.
\vfill
\newpage
\item Proof by \emph{Strong} Induction
\vspace{1in}
\item Proposition: Let $a_n$ be the sequence such that $a_0=1,$ $a_1=4$ and $a_n=3a_{n-1}-2a_{n-2}.$ Then, for every $n \in \N \cup \{0\},$ $a_n=3 \cdot 2^n-2.$
\vfill
\vfill

\item Proposition: Every integer $n >1$ has a prime factorization.
\vfill
(Bonus:) Any postage, $n$, can be made from 3-cent and 7-cent stamps provided $n\in \N$ and $n>12.$\end{enumerate}
\end{document}  

We proceed by induction on $n$.

\medskip\noindent(1) \textbf{Base case ($n=1$):}
The left-hand side equals $1\cdot 2 = 2$.
The right-hand side equals $\dfrac{1\cdot 2\cdot 3}{3} = 2$.
Hence the statement holds for $n=1$.

\medskip\noindent(2) \textbf{Inductive step:}
Assume that for some $k\geq 1$,
\[
  1\cdot 2+2\cdot 3+\cdots+k(k+1)=\frac{k(k+1)(k+2)}{3}.
\]
We must show the statement holds for $k+1$, i.e.,
\[
  1\cdot 2+2\cdot 3+\cdots+k(k+1)+(k+1)(k+2)=\frac{(k+1)(k+2)(k+3)}{3}.
\]
Starting from the left-hand side, observe
\hspace*{-.3in}\begin{align*}
  \big(1\cdot 2+\cdots+k(k+1)\big)+(k+1)(k+2)
  &= \frac{k(k+1)(k+2)}{3}+(k+1)(k+2)&\text{ by the inductive hypothesis}\\
  &= (k+1)(k+2)\!\left(\frac{k}{3}+1\right)& \text{ factoring}\\
  &= (k+1)(k+2)\cdot\frac{k+3}{3}& \text{ common denominator}\\
  &= \frac{(k+1)(k+2)(k+3)}{3}.&
\end{align*}
This is exactly the right-hand side for $n=k+1$, completing the induction.
\end{proof}














